\section{Related Work and Baselines}
\subsection{Clinical Feature-Based Cardiac Prediction}
Prediction on the UCI Cleveland dataset has moved from the logistic probability algorithm of Detrano \emph{et al.}~\cite{detrano1989} (about 77\% accuracy) to machine-learning and ensemble methods. Examples are the hybrid random forest with a linear model (HRFLM) of Mohan \emph{et al.}~\cite{mohan2019} (88.7\% accuracy), the majority-voting ensemble of Latha and Jeeva~\cite{latha2019} (85.48\%), the feature-selection study of Haq \emph{et al.}~\cite{haq2018} (89\% with logistic regression on Relief-selected features, 10-fold CV), and bagged decision trees with PCA by Gao \emph{et al.}~\cite{gao2021} (98.6\%). Results on the different Z-Alizadeh Sani dataset (303 patients, 54 features), such as the 93.85\% of Arabasadi \emph{et al.}~\cite{arabasadi2017}, are sometimes quoted alongside Cleveland results but concern another cohort. A limitation shared by Cleveland studies is the small cohort ($N=303$): a 10-fold test fold holds about 30 patients, so accuracy estimates vary considerably between folds and their confidence intervals are wide. Reported accuracies near or above 95\% should therefore be compared with others only when the validation protocol, in particular whether feature extraction and resampling were fitted inside the folds, is known to be equivalent. Section~VI-D shows that one common protocol error, oversampling before cross-validation, inflates AUC on this dataset.

\subsection{ECG Signal Analysis}
PTB-XL~\cite{wagner2020} enabled systematic benchmarking of ECG classifiers. Hannun \emph{et al.}~\cite{hannun2019} achieved cardiologist-level classification of 12 rhythm classes from single-lead ambulatory ECGs with a 34-layer deep neural network trained on more than 90{,}000 recordings. Ribeiro \emph{et al.}~\cite{ribeiro2020} trained a deep residual network on about 2.3 million 12-lead ECGs and reported F1 scores above 80\% for six abnormalities, outperforming cardiology residents. These studies differ from ours in two important ways: they use far larger training sets, and they predict \emph{ECG diagnostic statements} or rhythms, not angiographically confirmed coronary artery disease. Their accuracy figures are therefore not comparable with Cleveland results.

\subsection{Multimodal and Cross-Modal Fusion}
Combining data types is a major direction in medical AI, and a systematic review of imaging--EHR fusion~\cite{huang2020fusion} describes early (feature-level), joint and late fusion strategies. A single modality can also carry information usually obtained elsewhere: Attia \emph{et al.}~\cite{attia2019} detected left-ventricular systolic dysfunction (ejection fraction $\leq$35\%) from the 12-lead ECG alone with AUC 0.93. Cross-attention~\cite{vaswani2017} is increasingly used to let one modality condition on another. \emph{True} multimodal fusion, however, requires that the modalities describe the same patient, visit or event~\cite{huang2020fusion}. When they do not, as with Cleveland and PTB-XL, a model can at best approximate one modality from the other (proxy fusion) or evaluate each modality separately (parallel benchmarking). MultiCardioFusion uses proxy-aligned fusion on Cleveland and parallel ECG benchmarking on PTB-XL (Section~I).

\subsection{Explainability in Clinical AI}
SHAP~\cite{lundberg2017} is widely used for clinical models on structured data, and TreeSHAP~\cite{lundberg2018tree} computes exact Shapley values for tree ensembles such as XGBoost efficiently. In clinical use, Lundberg \emph{et al.}~\cite{lundberg2018sepsis} showed anaesthesiologists real-time SHAP explanations of intra-operative hypoxaemia risk. Three caveats apply. First, SHAP values are \emph{post-hoc and associational}: they describe the model, not causal mechanisms in the patient~\cite{kumar2020problems}. Second, with small samples and correlated or duplicated features, attributions can be split arbitrarily between features and can change between training runs~\cite{molnar2022pitfalls}. Third, agreement with known risk factors is necessary but not sufficient evidence of validity, and attributions can be manipulated~\cite{slack2020fooling}. We therefore treat SHAP as a descriptive tool, quantify its stability (Section~VI-H), and propose clinician review as future validation.

\begin{table*}[!t]
\centering\scriptsize
\caption{Literature summary and positioning of this work. Reported results are quoted as published; datasets, validation protocols and metrics differ between studies, so the figures are not directly comparable.}
\label{tab:lit}
\setlength{\tabcolsep}{3pt}
\begin{adjustbox}{max width=\linewidth}
\begin{tabular}{p{2.6cm}p{2.4cm}p{3.2cm}p{3.2cm}p{5.2cm}}
\toprule
Study & Data & Technique & Reported result & Relation to this work\\
\midrule
Detrano \emph{et al.}~\cite{detrano1989} 1989 & Cleveland (303) + other sites & Logistic probability algorithm & Acc.\ $\approx$77\% & Origin of the dataset and of the CAD label.\\
Haq \emph{et al.}~\cite{haq2018} 2018 & Cleveland (303) & 7 classifiers, Relief/mRMR/LASSO selection & Acc.\ 89\% (LR + Relief, 10-fold CV) & Clinical features only; no uncertainty estimates.\\
Mohan \emph{et al.}~\cite{mohan2019} 2019 & Cleveland (303) & HRFLM (RF + linear model) & Acc.\ 88.7\% & Clinical features only.\\
Latha \& Jeeva~\cite{latha2019} 2019 & Cleveland (303) & Bagging, boosting, stacking, majority voting & Acc.\ 85.48\% (majority voting) & Clinical features only; no explanations.\\
Gao \emph{et al.}~\cite{gao2021} 2021 & Cleveland (303) & Bagging/boosting with PCA/LDA & Acc.\ 98.6\% & Comparability depends on protocol details (Section~VI-D).\\
Arabasadi \emph{et al.}~\cite{arabasadi2017} 2017 & Z-Alizadeh Sani (303, 54 features) & GA-initialised neural network & Acc.\ 93.85\% & Different dataset; listed because it is often compared with Cleveland.\\
Hannun \emph{et al.}~\cite{hannun2019} 2019 & Single-lead ambulatory ECG ($>$90k) & 34-layer deep network & Cardiologist-level F1 (rhythm) & Rhythm labels; no clinical variables.\\
Ribeiro \emph{et al.}~\cite{ribeiro2020} 2020 & 12-lead ECG (2.3M) & Deep residual network & F1 $>$80\%, 6 abnormalities & ECG diagnostic labels, not CAD.\\
Attia \emph{et al.}~\cite{attia2019} 2019 & 12-lead ECG & Convolutional network & AUC 0.93 (EF $\leq$35\%) & ECG-only; different task.\\
\midrule
\textbf{This work (v5)} & Cleveland + 3 external sites; PTB-XL (parallel) & Proxy-enhanced clinical features; soft voting; dual-branch DL; SHAP & AUC \fullAUC{} \fullAUCci{} (5$\times$10 CV); clinical-only \clinAUC{} & Proxy fusion and parallel ECG benchmarking (not paired fusion); full ablation with CIs, controls and external validation.\\
\bottomrule
\end{tabular}
\end{adjustbox}
\end{table*}
